% Chapter 7: Using the library (practice author). PLAN.md section 1.4; revised for round 1 (responses/R1-07-using.md).
% Floats: T1 tab:method-choice (table*), F16 fig:newton and F18 fig:kink (margin); sidebars S13, S14.
% Examples: E11, E12, E13, E14, E15, E18, E19, E20, E22, E23 (code/07-using-*.cpp); the CMake file of 7.1.
\chapter{Using the library}\label{ch:using}

This chapter is for the reader with a problem: you have a function and a point, and you need a derivative. It shows
how to get the library and call it, and how to read and check the result. It then helps you choose a method, and works
through three engineering examples. The last sections cover other number types, use at compile time, and the
pitfalls. \Cref{ch:api} is the exact reference, and \cref{ch:tables} holds the tables.

\section{The call}\label{sec:using-call}

\paragraph{Getting the library.} \nxx{} is header-only and has no dependencies. It ships no install rule and no
\code{find_package} configuration. So a CMake project adds it as a subdirectory, with \code{add_subdirectory}, CPM or
\code{FetchContent}. The target \code{numerixx::numerixx} adds the include path and requires C++23. Here it is with
\code{FetchContent}, at a pinned commit or tag:
\cmakefile{getting-started}
The library's tests and this book's examples run in the continuous integration\mnote{\emph{Continuous integration}:
a service that builds and tests the code automatically after every change.} on GCC~16, on Clang~22 with libc++, on
MSVC and clang-cl of Visual Studio~2026, and on Emscripten~6.0.10. These are the toolchains that are tested. They are
not minimum versions.

\paragraph{One function.} Everything is in the header \code{<numerixx/deriv.hpp>}, in the namespace
\code{nxx::deriv}. There is one function, \code{derivative}, and it comes in two forms:
\begin{itemize}
  \item \code{derivative<M, N>(fn, x)} uses the method's default step;
  \item \code{derivative<M, N>(fn, x, h)} uses the step \code{h} that you pass (for the Ridders methods, the initial
    step).
\end{itemize}
The method \code{M} is a template argument. When you leave it out, it is \code{Central<2>}. Otherwise it is one of
\code{Central<P>}, \code{Forward<P>}, \code{Backward<P>}, \code{Ridders}, \code{ForwardRidders} and
\code{BackwardRidders}. The derivative order N is 1 by default. It is an integer from 1 to 4, or a value of the enumeration \code{order}, such as
\code{order::second}. The number type T is deduced from \code{x}: the compiler takes it from the type of \code{x}. A
step of another type converts to T, so \code{0.01} is a valid step for a 50-digit \code{x}. For a \code{float}
\code{x}, write the step as \code{1e-3f}: a \code{double} step is diagnosed under \code{-Wconversion} unless
\code{float} holds it exactly. \code{fn} is any callable\mnote{A \emph{callable} is anything \code{std::invoke} can
call: a function, a lambda, a function object, or a pointer to a const member function of T itself. For a member
function of your model, capture the object: \code{[&m](double t) { return m.H(t); }}.} that takes a T and returns
something convertible to T. The call of \cref{ch:why}, \code{derivative(H, T0)}, is the shortest form. It means
\code{Central<2>}, the first derivative and the default step. \Cref{sec:api-derivative} shows one call of each form,
compiled.

A stencil calls \code{fn} once per node. So \code{Central<2>} costs two evaluations for $f'$, and three for $f''$.
Ridders' method costs two or four evaluations per level, that is, for each step that it tries. In the random study
of \cref{ch:ridders}, its median was \num{\valRiddersEvalsMedianFirst} evaluations for $f'$ and
\num{\valRiddersEvalsMedianThird} for $f'''$ (\cref{tab:cost}).

\section{The result}\label{sec:using-result}

The library itself throws nothing: a failure is a value. An exception thrown by \code{fn}, or by the operations of T,
passes through unchanged, and \code{derivative} is not \code{noexcept}. The result holds either an estimate or the
reason why there is none. Its type is \code{std::expected<Estimate<T>, error>}. An \code{Estimate<T>} has two
members. \code{value} is the estimate of the derivative. \code{abs_error} is a \code{std::optional<T>}, which only
Ridders' method and its one-sided forms fill. So far there is one reason for a failure, \code{error::non_finite}.
Test the result, then read it:

\snippet{07-using-expected}{expected-basics}
The central rule at $x = 0$ samples $\log(-h)$, which is NaN, so the first call fails. \code{error::non_finite}
covers every case in which no finite estimate exists: \code{x} is not finite; \code{fn} returned an infinite or NaN
value; a difference, term, sum or quotient would overflow; a point $x + k h$ would overflow; or $h$ is below the
resolution of $x$, the spacing of the numbers near $x$ (\cref{sb:ulp}). \code{error} may gain values in a later
version, so a \code{switch} on it needs a \code{default}. The error is formattable:
\code{std::println("{}", d.error())} writes the name of its value, \code{non_finite}. To find the cause of a failure,
evaluate \code{fn} at $x \pm h$ yourself, and compare $h$ with the spacing of the numbers at $x$.

The member functions \code{transform} and \code{and_then} of \code{std::expected} apply a function only when there is
a value. A failure they pass through unchanged. \code{and_then} chains a second call, which can itself fail:
\needspace{13\baselineskip}% the listing in one piece (stage-6 layout: two comment lines were stranded)
\snippet{07-using-expected}{expected-chain}
Both lambdas take a \code{double}. That works because an \code{Estimate<double>} converts to its value wherever a
\code{double} is expected. The conversion is convenient, but it has traps, which the next snippet shows:
\snippet{07-using-expected}{estimate-pitfalls}
First, \code{auto} deduces \code{Estimate<double>}, and \code{std::println} does not apply the conversion. Print
\code{e.value} instead. Second, \code{if (*r)} tests the value against zero, not whether \code{r} holds a result. The
derivative of $\cos x$ at 0 is exactly 0: a good result, and it tests false. Clang with \code{-Wconversion -Werror}
rejects that line; GCC accepts it without a warning. Third, \code{r.value()} is the accessor of \code{std::expected},
not the member of the estimate. It compiles, converts to \code{double}, and works until the first failure. There it
throws \code{std::bad_expected_access}, or terminates a program built without exceptions. Read \code{r->value} after
the test. Fourth, \code{r.value_or(0.0)} does not compile, since no \code{double} converts to an
\code{Estimate<double>}; \code{r->abs_error.value_or(0.0)} does. And \code{int n = *r} truncates the value, with a
\code{-Wfloat-conversion} warning on GCC and Clang. The rest of the book reads a result as \code{d->value}, and
Ridders' estimate as \code{*r->abs_error}.

\paragraph{Did it work, and is it right?} \code{if (d)} answers the first question only. Each of these gives a wrong
value without a failure: a step on the wrong scale (\cref{sec:rounding-scale}), a function less accurate than
$\epsmach$ (\cref{sec:rounding-noise}), and a break within a step of $x$ (\cref{sec:using-pitfalls}).
\Cref{sec:using-check} shows how to check such a value. Ridders' method can also return a value whose own estimate
says that it is poor: near a bad point, or at the resolution of $x$. Before you use a Ridders result where it
matters, compare its estimate with the error that you can accept. The check uses two tolerances:
\needspace{15\baselineskip}% the listing in one piece (stage-6 layout: its last three lines were stranded)
\snippet{07-using-expected}{ridders-check}
For a value that you report to $n$ significant figures, \code{rtol} is $10^{-n}$. \code{atol} is the error that you can accept, in the derivative's units, and it is the one
that decides near a zero of $f'$ (\cref{sec:using-check}). Inside an iteration, no check is needed
(\cref{sec:using-newton}). A passed check says only that the estimate is small. It does not say that the value is
right where the estimate is blind, as at a kink (\cref{sec:using-pitfalls}).

\section{Choosing a method}\label{sec:using-choosing}

\Cref{tab:method-choice} is the decision table: find the row that describes your situation. Two rules cover most
cases. Inside an iteration, an error estimate would not be read. There, use a stencil with the default step:
\code{Central<2>} or \code{Central<4>}. Their median errors for $\mathrm{e}^x$ are \num{\valExpCentralTwoRelError}
and \num{\valExpCentralFourRelError}. For a value that you will report, use Ridders' method. It costs
\num{\valRiddersEvalsMedianFirst} evaluations for $f'$ (the median), against two to six for the stencils. In return
it gives about three more digits than \code{Central<2>} (\num{\valHeadlineRiddersFirstDigits} against
\num{\valHeadlineFirstDigits}). That is as many as \code{Central<4>} gives, and no more than \code{Central<6>}
(\cref{tab:accuracy-matrix}). And it gives an estimate of its error. The other rows of the table are the exceptions.

\begin{table*}
\caption{Choosing a method. The accuracies are relative errors in double. Where a cell names the model, it is the
model $\epsmach^{P/(P+N)}$ of \cref{ch:stencils}; the other values are measured. They are medians for
$\mathrm{e}^x$ (\cref{tab:accuracy-matrix}), or values for the nitrogen $\cp$ and this chapter's examples. The
evaluations are those for $f'$ unless the row says otherwise. For Ridders' method they are the study's median and the
cap (\cref{tab:cost}). $w$ is the scale on which $f$ varies, and $\noise$ is its relative accuracy. Above $f'$,
\code{float} warns (\cref{sec:api-warning}).\label{tab:method-choice}}
\centering\footnotesize
\begin{tabular}{@{}>{\raggedright}p{31mm}>{\raggedright}p{27mm}>{\raggedright}p{21mm}>{\raggedright}p{15mm}%
>{\raggedright}p{39mm}>{\raggedright\arraybackslash}p{19mm}@{}}
\toprule
Situation & Method & Step & Evaluations & Accuracy to expect & Evidence\\
\midrule
Inside an iteration; $f$ smooth and cheap & \code{Central<2>}, or \code{Central<4>} & default & 2, or 4 &
  model $\epsmach^{2/3} = \num{\valModelCentralTwoFirstDouble}$, measured median \num{\valExpCentralTwoRelError};
  \code{Central<4>} \num{\valExpCentralFourRelError} & \cref{fig:default-step-accuracy}, \cref{tab:accuracy-matrix}\\
A value to report, with an error estimate & \code{Ridders} & default initial step & \num{\valRiddersEvalsMedianFirst},
  at most \num{\valRiddersMaxEvalsFirst} & about \num{\valHeadlineRiddersFirstDigits} digits; the estimate covers the error
  in \SI{\valRiddersCoverageFirst}{\percent} of the study & \cref{fig:ridders-coverage}\\
$f$ defined on one side of $x$ only (an edge of its domain or its validity range) & \code{Forward<P>} or
  \code{Backward<P>}, P = 2 or 4; with an error estimate, \code{ForwardRidders} or \code{BackwardRidders} & default
  & 3, or 5; with the estimate \num{\valForwardRiddersEvalsMedianFirst}, at most
  \num{\valForwardRiddersMaxEvalsFirst} & \num{\valExpForwardTwoRelError} (P = 2); with the estimate
  \num{\valExpForwardRiddersFirstRelError}, whose estimate covers the error in
  \SI{\valForwardRiddersCoverageFirst}{\percent} of the study & \cref{fig:one-sided}, \cref{sec:stencils-one-sided,sec:ridders-onesided}\\
$f$ varies on a scale $w$ much smaller than $\abs{x}$, or $\abs{x} \ll 1$ in your units & \code{Central<2>}, or
  \code{Ridders} from $h = w/4$ & $\epsmach^{1/3} w$ & 2, or Ridders' \num{\valRiddersEvalsMedianFirst} &
  \num{\valScaleFitRelError}, against
  \num{\valScaleDefaultRelError} with the default step & \cref{fig:step-scale,fig:step-scale-small}\\
$f$ is the output of an iterative solver, accurate only to $\noise$ (noisy) & \code{Central<2>} & $\noise^{1/3}
  \max(\abs{x}, 1)$ & 2 & model $\noise^{2/3} = \num{\valSolverModelFitRelError}$, measured
  \num{\valSolverFitRelError} for $\noise = 10^{-10}$, against \num{\valSolverDefaultRelError} with the default step;
  Ridders' estimate is optimistic & \cref{fig:noise}, \cref{sec:using-solver}\\
A second derivative & \code{Central<2>} or \code{Ridders}, with $N = 2$; never float & default & 3, or
  \num{\valRiddersEvalsMedianSecond} & model $\epsmach^{1/2} = \num{\valModelCentralTwoSecondDouble}$, measured
  \num{\valCpSecondRelError}; Ridders \num{\valCpRiddersSecondRelError} (nitrogen) & \cref{fig:v-curve-orders}\\
A third or fourth derivative & \code{Ridders}, or \code{Central<4>}; never float & default & 4 per level, or 6 to 7 &
  $f''''$: Ridders \num{\valDoubleFourthRiddersRelError}, \code{Central<2>} \num{\valDoubleFourthCentralRelError};
  \code{Central<4>} in \cref{tab:accuracy-matrix} & \cref{fig:order-accuracy}\\
A wide type (50 digits) & \code{Central<6>} or \code{Ridders} & default & 6, or \num{\valMpRiddersEvals} &
  \num{\valMpCentralSixRelError}; Ridders \num{\valMpRiddersRelError} & \cref{sec:using-types}\\
A kink or a range break within a step of $x$ & none: differentiate one side, with \code{Forward<P>} or
  \code{Backward<P>}, or with an estimate \code{ForwardRidders} or \code{BackwardRidders} & default & as above & every
  method that samples across the break returns a wrong value; only Ridders' estimate shows a jump, and nothing shows
  a kink at $x$ &
  \cref{fig:kink}\\
\bottomrule
\end{tabular}
\end{table*}

\section{Worked examples}\label{sec:using-examples}

This section works through three problems from engineering practice: a temperature found by Newton's method, the
sensitivities of a rate constant together with the Jacobian of a set of balances, and the derivative of a function
that a solver computes.

\subsection{Newton's method for a temperature}\label{sec:using-newton}

An energy balance often gives the enthalpy of a stream and asks for its temperature. We must then solve
$\enth(T) = H^\ast$ for $T$. Newton's method needs the slope $\deriv{\enth}{T} = \cp$ at each iteration, and
\code{derivative} supplies it. This is the classic use of a derivative inside an iteration, so the method is the
default, \code{Central<2>}, applied to the nitrogen enthalpy \code{H(T)}. The solver returns two kinds of failure as
errors: a failed slope, and a failure to converge. So a caller cannot mistake an unconverged temperature for a
solution:
\snippet{07-using-newton}{newton}
Started at \SI{\valNewtonStart}{K} for the enthalpy of \SI{1000}{K}, it converges in \num{\valNewtonIterations}
iterations, to within \SI{\valNewtonFinalError}{K} (measured, \cref{fig:newton}). The final residual, the amount by
which $\enth(T)$ misses $H^\ast$, is \SI{\valNewtonFinalResidual}{\joule\per\mole}: the rounding of $\enth$ itself.

Why does a slope with an error do no harm? A slope that is off by a relative $\delta$ adds a term of about $\delta$
times the error to each step (Dennis and Schnabel,\cite{dennis1996} Chapter 5). Here $\delta$ is near
\num{\valCpDefaultRelError}, far too small to slow the convergence before rounding stops it. Each iteration costs
three evaluations of $\enth$: two for the slope, and one for the residual. Ridders' method would cost
\num{\valCpRiddersEvals} per slope, and buy nothing that the iteration can use.
% The margin figure is anchored after the paragraph that cites it so that it does not share a page with the
% full-width Table 7.1 (stage-7 layout check).
\begin{marginfigure}
\centering
% F16 fig:newton (Chapter 7, section 7.4): Newton's method for T with H(T) = H(1000 K) from 800 K, the slope by
% Central<2> with the default step (E12). The residual |H(T) - H*| in J/mol against the iteration: the number of
% correct digits doubles per step until the residual reaches the rounding of H. Data: DS11, data/newton.csv.
\begin{tikzpicture}
\begin{axis}[eqtufte, width=38mm, height=34mm, ymode=log,
  xlabel={iteration}, ylabel={residual $|H(T) - H^\ast|$ (J/mol)},
  xtick distance=1, ytick={1e-12,1e-8,1e-4,1,1e4},
  ymin=1e-13, ymax=1e5, xmin=-0.2, xmax=4.4]
  \addplot[accent, mark=*, mark size=1pt]
    table[col sep=comma, x=iteration, y expr={abs(\thisrow{residual})}, comment chars={\#}] {data/newton.csv}
    node[pos=0.5, anchor=south west, font=\scriptsize, accent, align=left] {quadratic:\protect\\digits double}
    node[pos=1, left=1pt, font=\scriptsize, muted] {rounding of $H$};
\end{axis}
\end{tikzpicture}

\caption{Newton's residual $\abs{\enth(T) - H^\ast}$ against the iteration, with the slope from \code{Central<2>}.
The number of correct digits doubles at each step, until the residual reaches the rounding of
$\enth$.\label{fig:newton}}
\end{marginfigure}

\subsection{Parameter sensitivities and a Jacobian}\label{sec:using-sensitivity}

A sensitivity is the derivative of a model output with respect to one input, with the other inputs held fixed. Take
an Arrhenius rate constant, $k = A\,\mathrm{e}^{-E_a/(\Rgas T)}$, with the illustrative values
$A = \SI{e10}{\per\second}$, $E_a = \SI{80}{\kilo\joule\per\mole}$ and $T = \SI{600}{K}$. We write it as a function
\code{rate(A, Ea, T)} of three arguments. \code{derivative} differentiates a function of one argument. A lambda that
captures the other two arguments makes one:
\snippet{07-using-sensitivity}{sensitivity}
These are values to report, so the method is Ridders'. The closed forms are $\partial k/\partial T = k
E_a/(\Rgas T^2)$ and $\partial k/\partial E_a = -k/(\Rgas T)$. Against them, the measured relative errors are
\num{\valSensitivityRateRelError} and \num{\valSensitivityEaRelError}, each below its estimate. Engineers usually
report the scaled sensitivity $\deriv{\ln k}{\ln E_a} = (E_a/k)\,\partial k/\partial E_a$ instead. It is the
relative change of $k$ per relative change of $E_a$, and it is dimensionless. Here it is \num{\valSensitivityScaled},
as is the exact $-E_a/(\Rgas T)$. So a small relative error in the activation energy reaches the rate magnified by
that factor. The sensitivities to all the parameters are the same call in a loop. Its lambda replaces one parameter
and keeps the rest:
\snippet{07-using-sensitivity}{jacobian}
The cost is the sum over the parameters: \num{\valJacobianEvals} evaluations in all. Of these,
\num{\valJacobianEvalsA} are for $A$. The rate is linear in $A$, so the first extrapolated entry agrees exactly with
its parent, and the method stops. And \num{\valJacobianEvalsEa} are for $E_a$.

\paragraph{The Jacobian of a set of balances.} Newton's method for several unknowns needs the Jacobian of several
residuals, and a model usually returns all its residuals from one call. Take an energy balance
$r_1 = n\,\enth(T) - Q$ and an ideal-gas pressure $r_2 = n \Rgas T / V - P$, in the unknowns $T$ and $n$. We work at
$T = \SI{1000}{K}$, $n = \SI{2}{mol}$ and $V = \SI{1}{\metre\cubed}$, with $Q$ and $P$ chosen so that both residuals
vanish there. The library differentiates a function of one argument with one value. So today it supports one way:
one \code{derivative} call per entry of the Jacobian, on a lambda that evaluates the whole model with one unknown
varied. That makes $m \times n$ calls, each costing two model evaluations with \code{Central<2>}. Here that is
\num{\valJacobianBalanceCallEvals} model evaluations, against \num{\valJacobianBalanceHandEvals} by hand (snippet
\code{jacobian-by-calls} in \code{code/07-using-jacobian.cpp}). For an expensive model, apply the central difference
of \cref{ch:taylor} to the whole vector of residuals yourself, at two model calls per unknown. The loop of
\code{jacobian(r, x)} in the same file does this; it fills \code{J[i][j]} $= \partial r_i / \partial x_j$:
\snippet{07-using-jacobian}{jacobian-balances}
Here that is \num{\valJacobianBalanceHandEvals} model calls, with steps of \SI{\valJacobianBalanceStepT}{K} and
\SI{\valJacobianBalanceStepN}{mol}. Every entry is within \num{\valJacobianBalanceHandMaxRelError} of the closed
forms $n \cp$, $\enth$, $n \Rgas / V$ and $\Rgas T / V$ (measured; the calls give
\num{\valJacobianBalanceCallMaxRelError}). The loop is the bare formula, without the library's checks. It does not
detect a non-finite residual, and it does not divide by the distance actually evaluated. So test the residuals that
it returns.

\subsection{A function with no formula: a friction factor}\label{sec:using-solver}

The nitrogen enthalpy has a closed-form derivative, but most functions in practice do not. The Darcy friction factor
$\lambda$ of turbulent flow in a rough pipe is an example. It follows from the Colebrook equation, which reads, in its
common form,\cite{colebrook1939}
\begin{equation}
  \frac{1}{\sqrt{\lambda}} = -2 \log_{10}\left(\frac{10^{-4}}{3.7} + \frac{2.51}{\mathit{Re}\,\sqrt{\lambda}}\right)
  \label{eq:using-colebrook}
\end{equation}
for a relative roughness of $10^{-4}$. This equation has no closed-form solution for $\lambda$. Instead,
$\lambda(\mathit{Re})$ is computed by an inner iteration: here, bisection on $1/\sqrt{\lambda}$ to a relative
tolerance $\noise$.\mnote{\emph{Bisection} finds a root of an equation by halving, again and again, an interval
that contains it.}
\snippetheader{colebrook}{implicit-function}
The output of such a solver is accurate to about its tolerance, not to machine epsilon. At $\mathit{Re} = 10^5$ with
$\noise = 10^{-10}$, it is \num{\valSolverFrictionNoise} off, relative. And it changes in steps as $\mathit{Re}$
moves. That is the noisy function of \cref{sec:rounding-noise}, and the cure from there applies: the step
$\noise^{1/3} s$ instead of the default $\epsmach^{1/3} s$.
\snippet{07-using-solver}{implicit-derivative}
The fitted step is the one to use, because the default step and Ridders' estimate both assume that $f$ is accurate to
$\epsmach$. The answer key is the derivative that follows from differentiating \cref{eq:using-colebrook}
implicitly, at its root, found to machine precision by Newton's method. With $h = \noise^{1/3}\,\mathit{Re}$, about
\num{\valSolverFitStep}, the relative error is \num{\valSolverFitRelError}. That is inside the model's
$\noise^{2/3} = \num{\valSolverModelFitRelError}$ (measured). With the default step, \num{\valSolverDefaultStep} in
units of $\mathit{Re}$, the error is \num{\valSolverDefaultRelError}. That is the rounding branch
$\noise \abs{\lambda} / h$ of the error, at a step much too small. Ridders' estimate happened to cover its error of
\num{\valSolverDerivRelError} here. But over \num{\valNoiseRiddersCases} random points of a function with relative
noise $10^{-8}$, it covered only \SI{\valNoiseEightRiddersCoverage}{\percent} of the cases
(\cref{sec:rounding-noise}). Here $\noise$ is known, because the solver's tolerance is on the output itself. When a
solver's test is on a residual instead, measure $\noise$ with the recipe of \cref{sec:rounding-noise}. On this
solver, the recipe returns \num{\valNoiseEstimateFriction} relative (the actual deviation is
\num{\valNoiseEstimateFrictionActual}), from \num{\valNoiseEstimateEvals} evaluations.

A solver can also be smooth but wrong by its tolerance. Then no method and no estimate can see the error. Ridders'
method on a fixed-point iteration\mnote{A \emph{fixed-point iteration} solves $x = g(x)$ by repeating
$x \leftarrow g(x)$.} stopped at the same tolerance was
\num{\valSolverFixedPointRiddersRelError} off, with a relative estimate of only
\num{\valSolverFixedPointRiddersRelEstimate} (measured). Converge the inner solver well below the accuracy that you
want of the derivative, or differentiate the converged equations implicitly.

\section{Is my derivative good enough?}\label{sec:using-check}

A stencil returns no error estimate. Four checks, in increasing order of cost, tell you whether to trust its value.
The first three use the nitrogen $\cp$ at $T_0 = \SI{1000}{K}$. They count the whole decimal digits on which two
values agree, that is, the integer part of $-\log_{10}(\abs{a - b} / \abs{b})$, with a helper \code{digits(a, b)}.

\paragraph{Repeat with $2h$ and $h/2$.} While truncation dominates the error, halving the step quarters the error of
\code{Central<2>}, whose truncation error is proportional to $h^2$. While rounding dominates, halving the step doubles
the error, since the rounding error is proportional to $1/h$. So values that agree to many digits at $h$, $2h$ and
$h/2$ are accurate to about that many digits. \code{default_step<M, N>(x)} returns the step that the call without a
step uses (\cref{sec:api-derivative}):
\needspace{4\baselineskip}% the lead-in with the listing (stage-6 layout); at 16 lines, after the language pass, the
% guard moved the listing off a page that had room for it and left half the page blank
\snippet{07-using-check}{check-by-steps}
Here the three values agree to \num{\valCheckStepsDigits} digits, and the value agrees with the closed form to
\num{\valCheckExactDigits} (measured). A jump inside the span, the interval that the stencil samples, does not survive
this check. A jump $J$ adds $J / (2h)$ to the central difference, and this term doubles as the step halves, so the
three values disagree. A kink at $x$ does survive the check: every step returns the mean of the two slopes
(\cref{sec:ridders-trust}).

\paragraph{Use a rule of higher order.} \code{Central<4>} has its own default step and a much smaller error. Here it
agrees with \code{Central<2>} to \num{\valCheckOrderDigits} digits. Since its own error is much smaller, this
agreement measures the error of \code{Central<2>}.


\paragraph{Ask Ridders' method.} It gives an independent value, with an estimate of its own error:
\snippet{07-using-check}{check-by-ridders}
Its estimate, \SI{\valCheckRiddersEstimate}{\joule\per\mole\per\kelvin}, belongs to Ridders' value, not to the
stencil's. But when the estimate is small, the difference
$\abs{d - r} = \SI{\valCheckRiddersDiff}{\joule\per\mole\per\kelvin}$ is the stencil's error, to within Ridders' own
error. Here Ridders' own error is \SI{\valCpRiddersActual}{\joule\per\mole\per\kelvin}, below its estimate. The
stencil's actual error is \SI{\valCheckStencilActual}{\joule\per\mole\per\kelvin}.

\paragraph{Where $f'$ is near zero, judge the absolute error.} Near a stationary point of $f$, where $f' = 0$, few
digits of $f'$ may agree although the derivative is accurate. The reason is that the error of the rule scales with
$\abs{f}$, not with $\abs{f'}$. Take $\sin x$ just above $\pi/2$. There \code{Central<2>} gives
$f' = \num{\valCheckSineDerivative}$, with an absolute error of only \num{\valCheckSineAbsError}, in units of $f$ per
unit $x$. Yet the three values at $h$, $2h$ and $h/2$ agree in only \num{\valCheckSineDigitsDerivative} digits of the
small slope (measured, \code{check-absolute} in \code{code/07-using-check.cpp}). So compare the spread of the values
with $\abs{f} / \max(\abs{x}, 1)$, not with $f'$.

Finally, a physical check costs nothing. A heat capacity is positive, and a reaction rate grows with temperature. A
derivative of the wrong sign is wrong, whatever the digits say.

\section{User-defined and multiprecision types}\label{sec:using-types}

\code{derivative} works on any type that satisfies the concept \code{nxx::real}.\mnote{A C++ \emph{concept} is a
named set of requirements on a type, which the compiler checks.} The concept includes the members of
\code{numeric_limits}, the standard description of a number type, that the library reads. The built-in
floating-point types satisfy it, and so do Boost.Multiprecision's numbers, or a type of your own
(\cref{sb:expression-templates}; \cref{sec:api-real} gives the list). The algorithms use only $+$, $-$, $\times$,
$/$, comparisons and \code{numeric_limits}. So a 50-digit type gets 50-digit derivatives:
\snippet{07-using-multiprecision}{multiprecision}
Take \code{cpp_bin_float_50}, whose machine epsilon is \num{\valEpsMultiprecision}. With it, \code{Central<6>} at its
default step is \num{\valMpCentralSixRelError} off, relative. Ridders' method is \num{\valMpRiddersRelError} off, with
a relative estimate of \num{\valMpRiddersRelEstimate}, after \num{\valMpRiddersEvals} evaluations (measured on
$\mathrm{e}^x$ at 1). \code{Central<2>} would stop at \num{\valMpCentralTwoRelError}. A wide type is where the higher
orders pay (\cref{tab:accuracy-matrix}). Two rules apply with every Boost type. First, in arithmetic, write
\code{mp50(*r)} or \code{r->value}, because Boost's operators do not see the conversion of an \code{Estimate}.
Second, give a lambda an explicit return type, \code{-> mp50}. The library converts a result at once, so an
expression over the parameter alone is safe. But an expression over the lambda's own locals would refer to destroyed
values (\cref{sb:expression-templates}). The explicit type is always safe.

\begin{sidebar}[label=sb:expression-templates]{Expression templates and user-defined real types}
\code{nxx::real<T>} asks for the following. \code{std::numeric_limits<T>} must be specialised as bounded, inexact and
non-integer, with the members that the algorithms read: \code{radix}, \code{digits}, \code{epsilon()}, \code{max()}
and \code{denorm_min()}. T must be constructible from \code{int}, copyable, and totally ordered, so that any two
values compare. And $+$, $-$, $\times$, $/$ and unary minus must have results that \emph{convert} to T. A type that
lacks one of these fails at the concept, at the call. \Cref{sec:api-real} quotes a minimal type that provides exactly
these members. Exact types such as rationals are rejected, since \code{numeric_limits} marks them exact (and
unbounded). Their epsilon is 0, and they have no largest value, so neither the step nor the overflow checks would mean
anything.

The results of the operators need not \emph{be} T. Boost.Multiprecision with expression templates (\code{et_on})
returns from \code{a + b} an expression object, not a number. The object refers to its operands, and it dangles once
they are gone: it then refers to values that no longer exist. The library never writes \code{auto} for an arithmetic
result, and it converts the value of \code{fn} with \code{static_cast<T>} at once. In user code, \code{-> T} on a
lambda does the same. Boost's arithmetic operators are templates that do not consider the conversion of an
\code{Estimate<T>}. So \code{abs(*r - e)} does not compile, with or without expression templates, while
\code{abs(T(*r) - e)} and \code{r->value - e} do.

\end{sidebar}

\section{At compile time}\label{sec:using-constexpr}

\code{derivative} is \code{constexpr}: the compiler can evaluate it.\mnote{A \code{constexpr} function can be
evaluated by the compiler, when its arguments are known at compile time.} When \code{fn} can itself be evaluated at
compile time, so can its derivative. The result can then feed a \code{static_assert}, a check that the compiler
makes, or a constant table. With a step that is a power of two, $x^3$ gives exact results in binary arithmetic. That
makes the truncation error of \cref{ch:taylor} visible to the compiler:
\snippet{07-using-constexpr}{constexpr}
Comparing the result with a value, as in \code{derivative(cube, 2.0, 0.25) == 12.0625}, compares \code{value} only,
through the conversion. The comparison is false when the result is a failure. But the same idiom with \code{!=} would
pass on a failure. Failures are values at compile time too. The nitrogen fit, written to return NaN below its
validity range, gives \code{error::non_finite} for \code{Central<2>} at \SI{500}{K}, and a value for
\code{Forward<2>}. Both results are checked by \code{static_assert} in \code{code/07-using-constexpr.cpp}. The one
exception is an overflow of the final quotient, for example in a fourth derivative with a tiny step. At run time it
gives \code{non_finite}, but GCC rejects it at compile time. Most of \code{<cmath>}, \code{exp} and \code{log} among
them, is not \code{constexpr} in C++23. So such an \code{fn} must be written with arithmetic only, as the nitrogen
enthalpy is. GCC evaluates some of these functions in constant expressions, as an extension, which Clang and MSVC
reject. So test compile-time derivatives on every toolchain that you ship.

The compiler and the processor compute the same value, provided that the compiler neither contracts
\code{a * b + c} into a fused multiply-add nor uses extended precision.\mnote{A \emph{fused multiply-add} computes
$ab + c$ with one rounding instead of two. \emph{Extended precision} keeps intermediate results with more digits than
the type has.} The build pins neither: it sets no option for them. The identity is verified on the five toolchains of
the continuous integration (\cref{sec:guarantees-guarantees}). How the library avoids the arithmetic that overflows is
the subject of \cref{sb:constexpr-overflow}.

\begin{sidebar}[label=sb:constexpr-overflow]{Overflow guards and constant evaluation}
At run time, an overflow gives an infinity, and a check afterwards finds it. In a constant expression it is worse.
GCC rejects any floating-point overflow and any NaN; Clang rejects a NaN and a division by zero. So a failure would
become a compile error instead of an error value. The library therefore detects each failure before the arithmetic
that would produce it. First, the input checks keep every point $x + k h$ at least about $K h$ below $\fmax$, the
largest finite value of the type, with $K$ the largest node; a NaN $x$ fails them too. Second, every difference and
partial sum is tested before it is formed. The test compares $a/\beta + b/\beta$ with $\fmax/\beta$, where $\beta$ is
the radix, the base in which the type stores its numbers. This comparison is exact and cannot overflow, so the NaN of
$\infty - \infty$ never arises. Third, the rounding floors of Ridders' method saturate at $\fmax$: they stop there
instead of overflowing. The exception is the final quotient, with its repeated divisions by $h$ for $\dorder \ge 2$.
Guarding it exactly would need overflow tests accurate to the last unit. So it is checked afterwards, which GCC
rejects in a constant expression.
\end{sidebar}

\section{Pitfalls}\label{sec:using-pitfalls}

Each pitfall below gives a wrong or missing result without a crash. For each one we give the symptom first, then the
cure.

\paragraph{The step's scale.} The default step assumes that $f$ varies on the scale $\max(\abs{x}, 1)$ in your units.
\emph{Symptom:} at a large $\abs{x}$ with a narrow feature, Ridders' method or a smaller step disagrees with the
default. At a small $\abs{x}$, such as a mole fraction of $10^{-6}$, the step is larger than $x$, and the result is
\code{non_finite} or crosses a physical bound. \emph{Cure:} pass $h = \epsmach^{1/3} w$, with $w$ the scale on which
$f$ varies (\cref{fig:step-scale,fig:step-scale-small}).

\paragraph{A kink or a range break inside the span.} A property model is often piecewise, that is, made of pieces:
two fits joined at a temperature, a phase boundary, an \code{if} in the code. The published enthalpy of nitrogen is
an example. It uses one Shomate fit below \SI{\valKinkJunction}{K} and another above, and the two do not meet:
$\enth$ jumps by \SI{\valKinkJumpEnthalpy}{\joule\per\mole} at \SI{\valKinkJunction}{K}.\cite{nist-webbook} Every
method evaluates it without complaint \SI{1}{mK} above the break:
\snippet{07-using-kink}{kink}
\begin{marginfigure}
\centering
% F18 fig:kink (Chapter 7, section 7.8): the published enthalpy of nitrogen across the range break at 500 K, where
% the two Shomate fits do not meet (H jumps by 0.77 J/mol). Relative error of cp at x = 500 K + delta against delta,
% default steps: Central<2>, Central<4>, Ridders' value and Ridders' relative error estimate (dashed). While the break
% lies inside a method's span the error is large and no method fails; only Ridders' estimate shows it. Data: DS20,
% data/kink.csv (E19). Labels sit on their curves (pos is a fraction of the drawn length), right of the drop of Central<2>.
\begin{tikzpicture}
\begin{axis}[eqtufte, width=38mm, height=56mm, xmode=log, ymode=log,
  xlabel={distance $\delta$ from the break (K)}, ylabel={relative error},
  xtick={1e-4,1e-2,1,1e2}, ytick={1e-14,1e-10,1e-6,1e-2,1e2},
  ymin=1e-15, ymax=1e2, xmin=4e-5, xmax=600, unbounded coords=jump]
  \addplot[muted]
    table[col sep=comma, x=delta, y=central2, comment chars={\#}] {data/kink.csv}
    node[pos=0, anchor=south west, font=\scriptsize] {\code{Central<2>}};
  \addplot[muted, densely dotted]
    table[col sep=comma, x=delta, y=central4, comment chars={\#}] {data/kink.csv}
    node[pos=0.15, anchor=south west, font=\scriptsize] {\code{Central<4>}};
  \addplot[black!80]
    table[col sep=comma, x=delta, y=ridders, comment chars={\#}] {data/kink.csv}
    node[pos=0.11, anchor=north west, font=\scriptsize] {Ridders};
  \addplot[accent, dashed]
    table[col sep=comma, x=delta, y=ridders_estimate, comment chars={\#}] {data/kink.csv}
    node[pos=0.11, anchor=south west, inner ysep=1pt, font=\scriptsize] {its estimate};
\end{axis}
\end{tikzpicture}

\caption{Relative error of $\cp$ against the distance $\delta$ from the range break at \SI{\valKinkJunction}{K}.
While the break lies inside a method's span, the error is a plateau. Its height is $J/h$ times the sum of the weights
on one side of the break, where $J$ is the jump in $\enth$. For \code{Central<2>} it is $J/(2h)$
(\num{\valKinkInsideCentralRelError} relative, $h = \SI{\valKinkCentralStep}{K}$). For \code{Central<4>} it is
$7J/(12h)$ while the break lies within $h$ of $x$ (\num{\valKinkInsideCentralFourRelError},
$h = \SI{\valKinkCentralFourStep}{K}$), and $-J/(12h)$ between $h$ and $2h$: the figure shows these two levels. So
the error falls as $1/h$ across the rules, whatever their order. No rule fails; only Ridders' estimate (dashed) shows
a jump.\label{fig:kink}}
\end{marginfigure}
All three calls succeed. \code{Central<2>} is \num{\valKinkInsideCentralRelError} off, relative: its error is the
jump divided by $2h$. \code{Central<4>} is \num{\valKinkInsideCentralFourRelError} off. So a higher order does not
cure a function that is not smooth. Ridders' value is \num{\valKinkInsideRiddersRelError} off. But its estimate,
\SI{\valKinkInsideRiddersEstimate}{\joule\per\mole\per\kelvin} or \num{\valKinkInsideRiddersRelEstimate} of the
value, is far above its usual size, which is \num{\valKinkOutsideRiddersRelEstimate} of the value away from the
break. The reason is that its extrapolation assumes an error expansion in even powers of $h$. The break destroys that
expansion, so the entries of the tableau disagree. At \SI{\valKinkOutsideDelta}{K} from the break, outside every
span, all three errors are at most \num{\valKinkOutsideCentralRelError} (measured, \cref{fig:kink}). Evaluable is not
smooth: that a function can be evaluated does not make it smooth. \emph{Symptom:} a large estimate from Ridders'
method, or stencil values that change with the step. A kink at $x$ shows no symptom at all, since central rules
return the mean of the two slopes (\cref{sec:ridders-trust}). \emph{Cure:} know the branch boundaries, and
differentiate the branch you mean. Above a break, use \code{Forward<P>}, or \code{ForwardRidders} for an error
estimate. Below it, use \code{Backward<P>} or \code{BackwardRidders}. Or call the branch's own function. At the break
itself, each of these returns the derivative of its own side, if $f(x)$ belongs to that side.

\paragraph{A Ridders result that is a value but poor.} Near a bad point, or at the resolution of $x$, Ridders'
method can return a value whose own estimate is a large fraction of it. \emph{Symptom:} \code{*r->abs_error} is far
above the accuracy that you need. \emph{Cure:} the check of \cref{sec:using-result}; a smaller initial step; or
\code{ForwardRidders} or \code{BackwardRidders} at an edge.


\paragraph{Aliasing.} A function that oscillates on a scale far below $\max(\abs{x}, 1)$ can lie about its slope
at every sample. One example is $\sin(64 \pi x)$ at $x = 0$, whose first levels sample only zeros
(\cref{sb:aliasing}). Another is $\sin x$ at a large $\abs{x}$, where the default step spans many periods.
\emph{Cure:} a step on the scale of the oscillation.

\paragraph{A noisy function.} A solver's output, or any function whose error changes irregularly from one argument
to the next, is accurate to $\noise$, not to $\epsmach$.
\emph{Symptom:} the error with the default step is far above $\epsmach^{2/3}$, and Ridders' estimate is optimistic.
\emph{Cure:} the step $\noise^{1/3} s$
(\cref{sec:using-solver}, \cref{sec:rounding-noise}).

\paragraph{A domain edge.} When $f$ is not finite at a step, Ridders' method restarts below that step. But it needs
the nearest bad point to be a minimum distance away. Take $\log x$ with the default step for $f'$. A value is
returned down to $x = 2^{-\num{\valRiddersEdgeLimitExponent}}$, with an estimate that has grown to flag the loss.
Ten correct digits hold only down to $x = 2^{-\num{\valRiddersEdgeTenDigitsExponent}}$ (measured,
\cref{fig:ridders-edge}). Closer to the edge, a NaN gives \code{non_finite}, and a pole, where $f$ grows without
bound, gives a flagged wrong value (\cref{sec:ridders-trust}). Exactly at an edge, use \code{ForwardRidders} or
\code{BackwardRidders} (\cref{sec:ridders-onesided}), or \code{Forward<P>} or \code{Backward<P>} inside an
iteration. For a fit evaluated outside its range, wrap the fit so that it returns NaN there, as
\cref{sec:stencils-one-sided} does. Then such a sample fails loudly.

\paragraph{Derivatives above the first in float.} They keep few correct digits, and every method gives a
compile-time warning (\cref{sec:orders-float,sec:api-warning}).
\emph{Cure:} use double.


\paragraph{One step for several orders.} The default step differs per method and order. With \code{Central<2>}, it
goes from $\epsmach^{1/3} s$ for $f'$ to $\epsmach^{1/6} s$ for $f''''$ (\cref{tab:default-steps}). It differs
because the rounding error grows as $\epsmach / h^N$. So a step tuned for $f'$ leaves nothing of $f''''$.
\emph{Cure:} in generic code that varies the order, pass no step.

\section{What the library does not do}\label{sec:using-not}

\code{nxx::deriv} differentiates a function, not a table. Tabulated data must first be fitted or interpolated, and
the derivative is then that of the fit. A linear interpolant, which joins the points by straight lines, has a
piecewise-constant derivative and a kink at every node. So start from a smooth fit or a spline,\mnote{A
\emph{spline} is a chain of polynomials, joined smoothly at the nodes.} and use the spline's own derivative when it
offers one. If you differentiate numerically, keep the stencil within one interval, or accept the error of a kink
(\cref{sec:using-pitfalls}).

\paragraph{Mixed partials and shared evaluations.} A mixed partial derivative such as
$\partial^2 A / \partial T \partial V$ takes two nested calls. The outer call is in $V$, on a lambda that calls the
inner one in $T$. The cost is the product of their costs. Let the lambda return the inner value, or a NaN when the
inner call fails (\code{d ? d->value : std::numeric_limits<double>::quiet_NaN()}). Then the outer call reports
\code{non_finite}. As a rule of thumb, the inner error, about $\epsmach^{2/3}$, is the noise level of the outer call.
The recipe for a noisy function then applies with $\noise = \epsmach^{2/3}$: pass the outer call the step
$\epsmach^{2/9} \max(\abs{V}, 1)$, and expect an error of about $\epsmach^{4/9}$, roughly $10^{-7}$ in double.
Finally, when an expensive call returns several outputs, cache its last argument and outputs, so that the outputs
share the evaluations.


